\section{编译器与框架}
\label{sec:compiler}

上一节的融合和分块，如果每个模型都要人手写，成本是不可接受的。框架决定了用户怎么写模型，编译器决定了硬件怎么跑它，
两者之间的那一层就是本节的主题。我们从计算图和自动微分讲起，接着讨论图级优化（融合、代数替换）、算子级优化（调度空间与搜索），
以及 PyTorch 2 在 Python 字节码层面捕获计算图的路线，最后专门讨论面试的高频难点：动态形状。

\subsection{从 define-and-run 到 define-by-run}

框架的第一个设计选择是：先构图再执行，还是边执行边构图。DistBelief~\cite{dean2012large} 和 TensorFlow~\cite{abadi2016tensorflow} 采用前者，
也就是静态图（define-and-run），好处是整张图可以统一优化和做分布式放置。PyTorch~\cite{paszke2019pytorch} 采用后者，即时执行（define-by-run），
直接用 Python 控制流表达模型，易用得多，代价是失去了全局优化的机会。JAX~\cite{jax2018github} 走的是函数变换的路线（\texttt{jit}、\texttt{grad}、\texttt{vmap}、\texttt{pmap}），
在 tracing 时得到纯函数的中间表示，再交给 XLA 编译~\cite{xla}。PyTorch 2~\cite{ansel2024pytorch2} 则想两者兼得：保留 eager 语义，在运行时把图捕获下来再编译。

\subsection{自动微分}

不管用哪种框架，训练都离不开自动微分。深度学习用的是反向模式自动微分~\cite{baydin2018automatic}：对标量损失 $\mathcal{L}$ 和 $n$ 个参数，
一次反传就能得到全部梯度，代价只是前向的常数倍；前向模式则需要 $n$ 次。对 GEMM 主导的网络，这个常数约为 2，
因为反向要对输入和对权重各做一次同样大小的 GEMM。训练 FLOPs 的经验公式 $C\approx6ND$ 就是这么来的：前向 $2ND$，反向 $4ND$。
反向模式的代价是必须保存前向的中间激活，训练显存里「激活」这一项就来源于此（第~\ref{sec:train} 节）。

\lab{15-849 编译与压缩}用 numpy 实现了一个反向模式自动微分，包括迭代式的拓扑排序、广播求和以及 softmax-交叉熵。
在两层 MLP 上，梯度与 \texttt{torch} 的最大差 $\le5.6\times10^{-17}$，与中心差分的最大差 $\le2.3\times10^{-10}$。
它还测试了 5000 层的链式乘法：因为拓扑排序是迭代而不是递归实现的，反向不会爆栈。

\begin{qa}{为什么深度学习用反向模式而不是前向模式自动微分？代价是什么？}
\ans 因为输出是一个标量，输入却是 $n$ 维参数。反向模式一次反传就拿到全部梯度，代价约为前向的 2--3 倍；前向模式要跑 $n$ 次（或用 $n$ 维的对偶数）。
代价是要保存前向激活，显存随层数和序列长度线性增长，这才引出了重计算和 offload。反过来，如果输出维度远大于输入，比如要求雅可比矩阵的列，前向模式更好。
\wrong 「反向模式更精确」。两者算出的都是精确导数，差别只在复杂度。
\end{qa}

\subsection{图级优化：融合与代数替换}

有了计算图，编译器就可以在图上做优化：常量折叠、公共子表达式消除、布局变换和融合。XLA 的 HLO 融合把逐元素算子和归约算子合成一个 kernel；
TVM~\cite{chen2018tvm} 的 Relay 层按算子类别（injective、reduction、complex-out-fusable）定下融合规则。
TASO~\cite{jia2019taso} 更进一步，自动生成并验证图替换规则，比如把共享同一输入的两个矩阵乘合成一个更宽的矩阵乘，再用代价驱动的回溯搜索选出替换序列。
融合并不总是划算。\lab{15-849 编译与压缩}的结果是，MLP 块的流量降了 49\%，时间却只降了 16\%（第~\ref{sec:kernel} 节）；
而且融合后寄存器压力变大，还可能降低 occupancy。

\subsection{算子级优化：调度空间与搜索}

图切好之后，每个算子内部的循环怎么排，是另一个问题。Halide~\cite{ragan2013halide} 提出了影响深远的「算法与调度分离」：
算法说明算什么，是一个纯函数定义；调度说明按什么顺序、在哪里、用几个核来算，用的是 \texttt{split}、\texttt{reorder}、\texttt{tile}、\texttt{vectorize}、\texttt{compute\_at} 这些原语。
改调度不会改变结果。TVM 把这个思想推广到张量程序，并用学习出来的代价模型指导搜索~\cite{chen2018tvm}。
Ansor~\cite{zheng2020ansor} 连调度模板也不需要人写了：它自动生成层次化的 sketch，再随机注解、做进化搜索，并用代价模型筛选候选。

\lab{15-849 编译与压缩}在 $2048^3$ 的 GEMM 上搭了一个调度空间，由循环顺序和 $t_m,t_n,t_k$ 组成；代价模型是解析的，
参数取单个 H100 SM 分到的算力和带宽（7.49~TFLOP/s、25.4~GB/s、228~KB）。在 3072 个调度里，只有 8 个落在最优值的 5\% 以内。
最优调度是 5.99~ms，也只达到峰值的 38\%，原因是 C 块的强度 $t_mt_n/(t_m+t_n)=128$ 低于拐点 295。
搜索到 5\% 以内需要评估多少次？随机搜索的中位数是 244、P90 是 649；进化搜索的中位数是 147.5、P90 是 496，大约快 1.65 倍；
按枚举顺序的网格搜索是 434，不过这个数只取决于枚举顺序，参考意义不大。

这个实验还给了我一个方法论上的教训。第一版代价模型误用了整颗 CPU 的参数，结果 1080 个调度都落在 5\% 以内，搜索问题被模型弄得过于简单。
可见代价模型决定了搜索问题有多难，这也是 Ansor 要依赖学习型代价模型和真实测量的原因。

\begin{qa}{TVM/Ansor 和 Triton 的定位有什么不同？}
\ans 两者解决的是不同的问题。TVM 把算法和调度分开，靠搜索为每个算子、每种硬件找调度（Ansor 自动生成 sketch、做进化搜索、用学习的代价模型），目标是「不用人写 kernel」。
Triton 让人写 block 级的程序，编译器负责 block 内的线程映射、访存合并和 shared memory，目标是「降低手写 kernel 的门槛」。
前者是自动调优框架，后者是一门编程语言加编译器。
\wrong 「Triton 是 TVM 的前端」；「Ansor 比 Triton 快」。两者解决的问题不同，没法直接比快慢。
\end{qa}

\subsection{PyTorch 2：TorchDynamo、AOTAutograd 与 TorchInductor}

回到「既要 eager 的易用，又要编译的性能」这个目标。PyTorch 2~\cite{ansel2024pytorch2} 的 \texttt{torch.compile} 由三部分组成。
TorchDynamo 借助 CPython 的帧求值钩子（PEP~523），在字节码层面捕获 FX 图，并为图生成 guard，也就是记录张量形状、dtype 和 Python 对象需要满足的条件。
AOTAutograd 在编译期追踪出前向和反向合在一起的联合图，这样反向也能被优化，并在这里决定哪些激活保存、哪些重算。
TorchInductor 把图降为 GPU 上的 Triton kernel，或者 CPU 上的 C++/OpenMP 代码。
遇到追踪不了的 Python 代码，比如数据相关的控制流、\texttt{.item()}、打印、不支持的第三方库，Dynamo 会发生 graph break：
断点前后变成两张独立的图，中间回到 Python 解释执行，跨越断点的融合机会也就没了。

\begin{qa}{\texttt{torch.compile} 什么时候会 graph break？代价是什么？怎么排查？}
\ans 会切图的情况有：数据相关的控制流（\texttt{if x.sum()>0}）、把张量转成 Python 标量（\texttt{.item()}、\texttt{.tolist()}）、打印和日志、
不支持的 C 扩展或第三方库。代价是每个断点都要回一次 Python，图被切碎以后跨断点的融合消失，还可能引起更多重编译。
排查可以用 \texttt{torch.\_dynamo.explain} 或 \texttt{TORCH\_LOGS=graph\_breaks}；修法是把控制流改写成 \texttt{torch.where} 或 \texttt{torch.cond}，
或者把追踪不了的部分移出编译区域。
\wrong 「compile 之后整个模型就是一张图」。
\end{qa}

\subsection{动态形状}

LLM 的序列长度和 batch 每一步都在变，编译器却偏爱静态形状，因为形状决定了分块方式、循环边界和内存规划。主流的处理策略有三种。
第一种是 padding 加分桶：把长度向上取整到若干个桶，用少量静态形状覆盖全部输入，这是 XLA/TPU 的传统做法，代价是 padding 上浪费的计算。
第二种是符号形状：PyTorch 2 把尺寸表示成符号（SymInt），在 guard 中记录形状之间的约束，只有约束被违反时才重编译~\cite{ansel2024pytorch2}。
第三种是运行时专用化：推理引擎为常见的 batch 大小预先捕获 CUDA Graph，其余情况走 eager 路径。
面试中要能讲清楚三者在编译次数、浪费的计算和 kernel 质量这三方面各自的取舍。

\begin{qa}{动态形状为什么让编译器难做？你会怎么处理变长 batch？}
\ans 因为形状决定了分块、循环边界、内存规划和 kernel 选择，形状一变，静态编译的产物要么失效需要重编译，要么变得次优。
训练时，我会按长度分桶并做打包（packing），减少 padding。推理时，decode 的 batch 维分桶并预先捕获 CUDA Graph；prefill 用 varlen attention，
按 cu\_seqlens 把序列拼接起来，避免 padding。在编译器一侧，用符号形状，并只把真正动态的维度标成动态（\texttt{mark\_dynamic}），免得每来一个新长度就重编译一次。
\wrong 「把所有输入 pad 到最大长度就行」。这会把 $T^2$ 的 attention 计算浪费在 padding 上。
\end{qa}

\begin{table*}[!t]
\centering
\caption{ML 编译器与编程系统对比}
\label{tab:compilers}
\footnotesize
\begin{tabularx}{\textwidth}{@{}l L{2.6cm} L{3.2cm} Y Y@{}}
\toprule
系统 & 输入/抽象 & 核心思想 & 优化手段 & 典型局限 \\
\midrule
Halide~\cite{ragan2013halide} & 图像流水线 DSL & 算法与调度分离 & 手写或自动调度（tile、vectorize、compute\_at） & 领域限于规则 stencil/张量 \\
TVM~\cite{chen2018tvm} & Relay 图 + 张量表达式 & 端到端张量编译，调度原语 & 图级融合 + 模板化调度 + 学习的代价模型 & 依赖调度模板；调优耗时 \\
Ansor~\cite{zheng2020ansor} & 张量表达式 & 无模板的层次化搜索空间 & sketch 生成、随机注解、进化搜索、代价模型 & 搜索时间长；对新硬件需要测量 \\
TASO~\cite{jia2019taso} & 算子图 & 自动生成并形式验证图替换 & 代价驱动的回溯搜索替换序列 & 主要是图级，不生成 kernel \\
XLA~\cite{xla} & HLO（JAX/TF） & 整图编译，静态形状 & 融合、布局、buffer 复用；GSPMD 自动分片 & 动态形状需 padding/分桶 \\
PyTorch 2~\cite{ansel2024pytorch2} & Python 字节码 & 保留 eager 语义的图捕获 & Dynamo guard、AOTAutograd、Inductor 生成 Triton & graph break；重编译 \\
Triton~\cite{tillet2019triton} & block 级 Python kernel & 以 tile 为单位编程 & 编译器负责线程映射、合并访存、shared memory & 需人写算法；跨 block 协作受限 \\
\bottomrule
\end{tabularx}
\end{table*}

把这些系统按时间放在一起看（表~\ref{tab:compilers}），能看到一条清楚的演化线：最早是人写调度（Halide、TVM 模板），然后是机器搜索调度（Ansor），
再到人写 tile 级算法、由编译器负责 tile 内部的映射（Triton），最后是框架自动捕获计算图，并把融合后的子图交给 Triton（PyTorch 2）。
GSPMD~\cite{xu2021gspmd} 和 Alpa~\cite{zheng2022alpa} 把同样的编译思想用到了分布式并行上，第~\ref{sec:train} 节会讲到。

\begin{keypts}
\begin{itemize}
\item 反向模式 AD：一次反传得到全部梯度，代价约为前向的 2 倍（$6ND$ 的来源），但需要保存激活。
\item 图级优化有融合、CSE 和代数替换（TASO）；算子级优化靠算法与调度分离（Halide/TVM）和搜索（Ansor）。
\item \texttt{torch.compile} 由 Dynamo（字节码捕获 + guard）、AOTAutograd（联合前反向图）和 Inductor（生成 Triton）组成；graph break 会切断融合。
\item 动态形状的三种策略：分桶 padding、符号形状、运行时专用化（CUDA Graph）。
\end{itemize}
\end{keypts}

\begin{research}
\begin{itemize}
\item 代价模型的可迁移性：在一类 GPU 上训练的代价模型，能不能迁移到新架构上？可以把本文的解析模型当作「先验」，
只用少量真实测量做校准，衡量搜索到 5\% 以内需要多少次测量。
\item 融合决策应该基于 roofline，而不是简单地「能融就融」。一个可做的研究是把每个候选融合组的收益估计为
$\Delta t=\max(F/P,B/\mathrm{BW})$ 的变化量，再与 Inductor 的启发式规则对比。
\end{itemize}
\end{research}
